% SYNC-ID: discussion
\section{讨论}\label{sec:discussion}
% CONTENT_DEFERRED_TO_WRITING_PHASE
\textit{(THIS IS PLACEHOLDER TEXT, PLZ FULFILL IT WHEN WRITING.)}
% SYNC-ID: discussion-placeholder-1
讨论将在论断、证据、限制和认识强度完成审核后撰写。本段只预留预期的阅读顺序。
% SYNC-ID: discussion-placeholder-2
后续将把观察与限制相连，但不在讨论中新增论断。
% SYNC-ID: discussion-placeholder-3
本段预留经过校准的解释，具体内容交由批准的证据记录决定。
% SYNC-ID: discussion-placeholder-4
复核后将明确标出仍待确认的表述。
% SYNC-ID: discussion-placeholder-5
范围比较应当简短、配对，并限于研究能够支持的内容。
% SYNC-ID: discussion-placeholder-6
正式文字将校准因果措辞，不暗示证据未支持的机制。
% SYNC-ID: discussion-placeholder-7
每个解释在接受本占位内容前都要连接到证据指针。
% SYNC-ID: discussion-placeholder-8
讨论的最后一段预留通向有效性与结论的克制过渡。正式版本应让观察、解释和开放问题之间的边界易于审计。
% SYNC-ID: discussion-placeholder-9
本 fixture 还预留一小段，用于说明当前记录无法证明的内容。措辞应和相应限制一起确定，不能在翻译时临时发挥。
% SYNC-ID: discussion-placeholder-10
讨论版式还留出少量空间，便于读者跟随从观察到受限解释的转折。正式版本应指出该转折的来源，并保留重要的不确定性。本通用句子只是排版 fixture，不声称当前存在任何观察、解释或开放问题。
